% problem-000036 6 钌配合物与钌红结构
\section*{6. 钌配合物与钌红结构}
\textit{下列为分问。}

\subsection*{6-1}
研究者制得⼀种含混合配体的Ru(II)配合物 $\mathrm { R u ( b p y ) } _ { n } ( \mathrm { p h e n } ) _ { 3 - n } ] ( \mathrm { C l O } _ { 4 } ) _ { 2 }$ （配体结构如下图）。元素分析结果给出C、H、N的质量分数分别为48.38%、3.06%、10.54%。磁性测量表明该配合物呈抗磁性。

\subsection*{6-1}
-1 推算配合物化学式中的�值。

\subsection*{6-1}
-2 写出中⼼钌原⼦的杂化轨道类型。

（图：）
bpy

（图：）
phen

\subsection*{6-2}
利⽤显微镜观察⽣物样品时，常⽤到⼀种被称为“钌红”的染⾊剂。钌红的化学式为

$\mathrm { [ R u _ { 3 } O _ { 2 } ( N H _ { 3 } ) _ { 1 4 } ] C l _ { 6 } }$ ，由 $\mathrm { [ R u ( N H _ { 3 } ) _ { 6 } ] C l _ { 3 } }$ 的氨⽔溶液暴露在空⽓中形成。钌红阳离⼦中三个钌原⼦均为6配位且⽆⾦属-⾦属键。

\subsection*{6-2}
-1 写出⽣成钌红阳离⼦的反应⽅程式。

\subsection*{6-2}
-2 画出钌红阳离⼦的结构式并标出每个钌的氧化态。

\subsection*{6-2}
-3 写出钌红阳离⼦中桥键原⼦的杂化轨道类型。

\subsection*{6-2}
-4 经测定，钌红阳离⼦中Ru—O键⻓为187pm，远⼩于其单键键⻓。对此，研究者解释为：在中⼼原⼦和桥键原⼦间形成了两套由�和�轨道重叠形成的多中⼼π键。画出多中⼼π键的原⼦轨道重叠⽰意图。
